\documentclass[11pt]{amsart}
\usepackage{amsmath,amssymb,amsthm}
\usepackage[margin=1.05in]{geometry}

\numberwithin{equation}{section}
\newtheorem{theorem}{Theorem}[section]
\newtheorem{proposition}[theorem]{Proposition}
\newtheorem{lemma}[theorem]{Lemma}
\newtheorem{corollary}[theorem]{Corollary}
\theoremstyle{definition}
\newtheorem{definition}[theorem]{Definition}
\theoremstyle{remark}
\newtheorem{remark}[theorem]{Remark}
\newtheorem{example}[theorem]{Example}

\newcommand{\Z}{\mathbb{Z}}
\newcommand{\Q}{\mathbb{Q}}
\newcommand{\R}{\mathbb{R}}
\newcommand{\N}{\mathbb{N}}
\newcommand{\C}{\mathbb{C}}
\newcommand{\hZ}{\widehat{\Z}}
\newcommand{\cA}{\mathcal{A}}
\newcommand{\cL}{\mathcal{L}}
\newcommand{\cO}{\mathcal{O}}
\newcommand{\Idl}{\mathbb{I}}
\newcommand{\pp}{\mathfrak{p}}
\newcommand{\aaa}{\mathfrak{a}}
\newcommand{\qq}{\mathfrak{q}}
\newcommand{\Ns}{\mathrm{N}}
\newcommand{\Cl}{\operatorname{Cl}}
\newcommand{\Lk}{\operatorname{Lk}}
\newcommand{\Vol}{\operatorname{Vol}}
\newcommand{\Gal}{\operatorname{Gal}}
\newcommand{\Frob}{\operatorname{Frob}}
\newcommand{\Prob}{\operatorname{Prob}}
\newcommand{\Spec}{\operatorname{Spec}}
\newcommand{\tlam}{\widetilde\lambda}

\PassOptionsToPackage{hyphens}{url}
\usepackage[hidelinks,bookmarks=false]{hyperref}

\title[Bost--Connes systems in arithmetic topology]{Localized Bost--Connes systems in arithmetic topology:\\ local class field theory, linking Kakutani dichotomies,\\ and topological phase transitions}
\author{Ruqing Chen}
\address{GUT Geoservice Inc., Montreal, Canada}
\email{ruqing@hotmail.com}
\thanks{Sources, code and data for the entire series: \url{https://github.com/Ruqing1963/localized-bost-connes}; this paper is in \texttt{papers/VIII-arithmetic-topology/}.}
\date{}

\begin{document}

\begin{abstract}
We quantize the Mazur--Kapranov--Reznikov dictionary \cite{Mazur,Kapranov,Reznikov,Morishita}
by transporting the localized Bost--Connes theory of \cite{Main} to knots in a
$3$-manifold, and determine which topological data drive the resulting thermodynamics.

Everything rests on one choice. In the dictionary the boundary torus of a knot plays the
role of the local field, $H_1(\partial N(K_i);\hZ)\leftrightarrow\widehat{K_\pp^\times}=\pi_\pp^{\hZ}\times\cO_\pp^\times$,
with the meridian corresponding to \emph{inertia} and the longitude to \emph{Frobenius}, so
the Bost--Connes parameter must be $\sigma_{K_i}=[\lambda_i]$. The half-lives-half-dies
theorem then completes the dictionary: the Lagrangian $\ker(H_1(\partial X)\to H_1(X))$,
spanned for $M=S^3$ by the $0$-framed longitudes $\tlam_i=\lambda_i-\sum_j\Lk(K_i,K_j)\mu_j$,
corresponds to the principal ideles $\overline{K^\times_{S,+}}$, its intersection with the
meridian subgroup to the global units $\overline{\cO^\times_{K,+}}$ --- which vanishes for
$M=S^3$, matching $\Z^\times_+=\{1\}$ under $\Spec\Z\leftrightarrow S^3$ --- and the subgroup
of $H_1(M;\Z)$ generated by the knot classes to $\Cl^+_S$.

With this parameter the $\mathrm{KMS}_\beta$ simplex is $\Prob(G_\cL/\Xi_\beta^\perp)$ with
the obstruction group governed by the Gauss linking matrix: a character with finite support
$F$ and denominator $m$ satisfies
$\chi(\sigma_{K_i})=\exp\bigl(2\pi i\sum_{j\in F}\Lk(K_i,K_j)a_j\bigr)$, so its detecting set
is $D_\chi=\{i:\sum_{j\in F}\Lk(K_i,K_j)b_j\not\equiv0\bmod m\}$. Two structural facts decide
whether $D_\chi$ is infinite, and we prove both: local finiteness is vacuous for an infinite
link in a closed manifold, so components must accumulate; and once they do, a compact Seifert
surface can be pierced infinitely often and the linking rows acquire infinite support. We
exhibit the simplest such family --- infinitely many parallel copies of a knot with framing
$f$ --- for which $D_\chi$ omits only $F$ whenever $f\sum_{j\in F}a_j\notin\Z$, so that
$\Xi_\beta$ is genuinely $\beta$-dependent, although for this family some characters are
never detected and the state is never unique. We then prove that every symmetric integer
matrix is the linking matrix of a countable tame link in $S^3$, and use the matrix of
binary digits $\Lk(K_i,K_j)=\mathrm{bit}_{\min(i,j)}(\max(i,j))$ to build a link whose
$\mathrm{KMS}_\beta$ state is unique for $\beta\le\beta_c$ and whose simplex is all of
$\Prob(\prod_j\hZ)$ above: the Bost--Connes phase diagram, realized by linking alone, with
open or closed broken phase according to the norm. Ambient homology enters twice: the
subgroup of $H_1(M;\Z)$ generated by the knot classes carries the exact analogue of the
$h^+_S>1$ obstruction, while the quotient $H_1(M;\Z)/\langle[K_i]\rangle$, which the
topological group sees and $G_S$ never does, breaks the symmetry at every temperature.

Section~\ref{sec:autopsy} isolates why the meridian is the wrong choice: $\mu_i$ is a single
basis vector of $H_1(S^3\setminus\cL)$, so $D_\chi$ is finite and $\Xi_\beta$ is everything at
every temperature. This is what inertia should do --- inertia is what an unramified cover
kills, and the phase structure lives in the unramified direction. The contrast with
$\sigma_\pp=(\pi_\pp)_{\qq\ne\pp}$, nontrivial at every other place by the product
formula, is exact: reciprocity and linking are two ways of spreading a local generator over
infinitely many places.
\end{abstract}

\maketitle

\tableofcontents

\section{The local dictionary and its homological foundation}\label{sec:local}

Let $M$ be a closed, connected, oriented $3$-manifold and $\cL=\{K_1,K_2,\dots\}$ a countable
family of disjoint tame knots. Write $N(\cL)$ for an open tubular neighbourhood,
\[
X:=M\setminus N(\cL),\qquad \partial X=\bigsqcup_iT_i,\qquad T_i=\partial N(K_i)\cong T^2,
\]
so $H_1(T_i;\Z)=\Z\mu_i\oplus\Z\lambda_i$ is free on the meridian and a chosen longitude.
For $M=S^3$ we take $\lambda_i$ to be the Seifert longitude, and we write $\Lk(K_i,K_j)$ for
the linking number when $i\ne j$ and $\Lk(K_i,K_i)$ for the framing of $\lambda_i$, so
$\Lk(K_i,K_i)=0$ in $S^3$.

\subsection{Local class field theory for knots}

\begin{definition}[The local entry]\label{lem:localcft}
The dictionary of Mazur, Kapranov--Smirnov, Reznikov and Morishita
\cite{Mazur,Kapranov,Reznikov,Morishita} pairs $\Spec\cO_\pp\leftrightarrow N(K_i)$ and
$\Spec K_\pp\leftrightarrow N(K_i)\setminus K_i\simeq T_i$, and the local reciprocity
isomorphism $\Gal(K_\pp^{ab}/K_\pp)\cong\widehat{K_\pp^\times}=\pi_\pp^{\hZ}\times\cO_\pp^\times$
\cite[Ch.~V]{Neukirch} with $\pi_1^{ab}(T_i)^\wedge=\hZ\lambda_i\oplus\hZ\mu_i$, the entries
being
\begin{equation}\label{eq:localdict}
\mu_i\ \longleftrightarrow\ \text{inertia }\cO_\pp^\times,
\qquad
\lambda_i\ \longleftrightarrow\ \text{Frobenius }\pi_\pp^{\hZ}.
\end{equation}
This is a convention, not a theorem, and everything below depends on it; the reasons for it
are the following. In a cyclic cover of $M$ branched along $K_i$ the meridian carries the
local monodromy around the branch locus, and the meridians are exactly what must die in an
unramified cover of $M$, as the inertia subgroup $\cO_\pp^\times$ is exactly what dies in an
unramified extension. The unramified quotient $\Gal(K_\pp^{ur}/K_\pp)=\hZ$ is generated by
Frobenius and matches $\hZ\lambda_i$, the longitude surviving in an unramified cover with its
image recording the covering degree along $K_i$. (The tame quotient of $\Gal(\bar K_\pp/K_\pp)$
is $\hZ'\rtimes\hZ$ and matches $\pi_1(T_i)$ only after abelianization and completion; the
matching of the inertia group $\cO_\pp^\times$ with the single generator $\mu_i$ is the
usual coarsening of the dictionary.)
\end{definition}

\begin{remark}[The choice that decides everything]\label{rem:choice}
The informal statement of the dictionary --- ``primes correspond to knots'' --- does not say
which class on the boundary torus is Frobenius, and the meridians are the visible generators
of $H_1(S^3\setminus\cL)$ (Alexander duality). Setting $\sigma_{K_i}=[\mu_i]$ is therefore tempting, and yields a
theory with no thermodynamics at all; \S\ref{sec:autopsy} carries out that computation and
explains it. Everything below uses $\sigma_{K_i}=[\lambda_i]$.
\end{remark}

\subsection{Half lives, half dies, and the unit group}

Let $\iota_*:H_1(\partial X;\hZ)\to H_1(X;\hZ)$ be induced by inclusion and write
$\Idl_\partial:=H_1(\partial X;\hZ)=\bigoplus_i(\hZ\mu_i\oplus\hZ\lambda_i)$.

\begin{theorem}[The Lagrangian]\label{thm:hlhd}
For $X$ compact oriented with boundary a disjoint union of $n$ tori,
\[
\dim_\Q\ker\bigl(H_1(\partial X;\Q)\to H_1(X;\Q)\bigr)=\tfrac12\dim_\Q H_1(\partial X;\Q)=n,
\]
and the kernel is Lagrangian for the intersection form. For $M=S^3$ it is spanned by the
$0$-framed longitudes
\begin{equation}\label{eq:lagr}
\tlam_i:=\lambda_i-\sum_j\Lk(K_i,K_j)\,\mu_j .
\end{equation}
\end{theorem}

\begin{proof}
The first statement is the classical half-lives-half-dies lemma, from the long exact sequence
of $(X,\partial X)$ with Poincaré--Lefschetz duality \cite[Ch.~4]{Thurston}. For $M=S^3$, Alexander duality gives
$H_1(X)=\bigoplus_j\Z\mu_j$ with $\iota_*\mu_i=\mu_i$ and
$\iota_*\lambda_i=\sum_j\Lk(K_i,K_j)\mu_j$, so $\sum_i(p_i\mu_i+q_i\lambda_i)$ lies in the
kernel precisely when $p_j=-\sum_iq_i\Lk(K_i,K_j)$ for every $j$; this is free of rank $n$ on
the $q_i$, with basis \eqref{eq:lagr}.
\end{proof}

\begin{theorem}[The completed dictionary]\label{thm:dict}
Let $\cL$ be finite with $n$ components, $\iota_*:\Idl_\partial\to H_1(X;\hZ)$, and let
$\nu:H_1(X;\hZ)\to H_1(M;\hZ)$ be induced by inclusion, so that $\nu\iota_*\mu_i=0$ and
$\nu\iota_*\lambda_i=[K_i]$. Put
\[
G_\cL^{\mathrm{idl}}:=\Idl_\partial/\ker\iota_*=\iota_*(\Idl_\partial)\le H_1(X;\hZ),\qquad
\Cl_\cL:=\bigl\langle[K_i]:i\bigr\rangle\le H_1(M;\hZ).
\]
Then the idele exact sequence
\[
1\longrightarrow U_S/\overline{\cO^\times_{K,+}}\longrightarrow G_S=\Idl_S/\overline{K^\times_{S,+}}\longrightarrow\Cl^+_S\longrightarrow1
\]
of \cite[Prop.~2.5]{Main}, in which $\Cl^+_S\le\Cl^+_K$ is the \emph{subgroup} generated by
the classes of the primes of $S$ and the last map is the divisor map, corresponds term by
term to the exact sequence
\[
1\longrightarrow\Bigl(\bigoplus_i\hZ\mu_i\Bigr)\Big/\Bigl(\ker\iota_*\cap\bigoplus_i\hZ\mu_i\Bigr)
\longrightarrow G_\cL^{\mathrm{idl}}
\xrightarrow{\ \nu\ }\Cl_\cL\longrightarrow1,
\]
via
\[
\Idl_S\leftrightarrow\Idl_\partial,\quad
U_S\leftrightarrow\bigoplus_i\hZ\mu_i,\quad
\overline{K^\times_{S,+}}\leftrightarrow\ker\iota_*,\quad
\overline{\cO^\times_{K,+}}\leftrightarrow\ker\iota_*\cap\bigoplus_i\hZ\mu_i,\quad
\Cl^+_S\leftrightarrow\Cl_\cL .
\]
For $M=S^3$ the unit term and $\Cl_\cL$ both vanish, matching $\Z^\times_+=\{1\}$ and
$\Cl(\Z)=1$ under $\Spec\Z\leftrightarrow S^3$, and $G_\cL^{\mathrm{idl}}=H_1(X;\hZ)=\bigoplus_i\hZ\mu_i$.
In general $H_1(X;\hZ)/G_\cL^{\mathrm{idl}}\cong H_1(M;\hZ)/\Cl_\cL$, a quotient with no
arithmetic counterpart.
\end{theorem}

\begin{proof}
The first two assignments are Definition~\ref{lem:localcft}, $K_\pp^\times$ decomposing into
a valuation part and a unit part matching $\hZ\lambda_i$ and $\hZ\mu_i$. The principal ideles
are those dying in the global quotient, matching $\ker\iota_*$; the global units are the
principal ideles that are local units at every place, matching the intersection with the
meridian subgroup; and the divisor map, which reads off the valuations of an idele, matches
$\nu$, which reads off the longitude coefficients since $\nu\iota_*\mu_i=0$ and
$\nu\iota_*\lambda_i=[K_i]$, the longitude being isotopic to the core of $N(K_i)$. Exactness:
$\nu$ maps $G_\cL^{\mathrm{idl}}$ onto the subgroup generated by the $[K_i]$; its kernel is
$\iota_*^{-1}(\ker\nu)/\ker\iota_*$, and $\ker\nu=\langle\iota_*\mu_i\rangle$ because $M$ is
obtained from $X$ by Dehn filling along the meridians, so that $H_1(M)=H_1(X)/\langle\mu_i\rangle$;
if $\iota_*\xi=\iota_*\eta$ with $\eta\in\bigoplus\hZ\mu_i$ then $\xi\equiv\eta$ modulo
$\ker\iota_*$, which identifies the kernel with the image of $\bigoplus\hZ\mu_i$. For
$M=S^3$, by \eqref{eq:lagr} an element of $\ker\iota_*$ is $\sum_iq_i\tlam_i$, which lies in
$\bigoplus\hZ\mu_i$ only if every $q_i=0$, and $H_1(S^3)=0$. The last assertion is
$H_1(X)/\langle\iota_*\mu_i,\iota_*\lambda_i\rangle=(H_1(X)/\langle\mu_i\rangle)/\langle[K_i]\rangle$.
\end{proof}

\begin{remark}[Numerical check]\label{rem:numlagr}
For random symmetric linking matrices with zero diagonal and $n\in\{2,3,5,8\}$, the rank of
$\ker\iota_*$ was computed to be exactly $n=\frac12\dim H_1(\partial X)$ in every case, and
$\ker\iota_*\cap\bigoplus\hZ\mu_i=0$ in every case.
\end{remark}

\section{The knot-theoretic Bost--Connes system}\label{sec:system}

\begin{remark}[Infinite links must accumulate]\label{rem:accumulate}
For infinite $\cL$ in a closed $M$ one cannot require that every compact set meet only
finitely many components: taking the compact set to be $M$ would force $\cL$ finite. The
components of an infinite link therefore accumulate on a nonempty closed set. This is not a
defect; by Theorem~\ref{thm:infinite} it is exactly what allows the linking rows to have
infinite support, and hence what makes the theory nonempty.
\end{remark}

\begin{definition}[Norm]\label{def:norm}
A norm is a function $N:\cL\to[2,\infty)$, extended to $J_\cL=\bigoplus_i\N\cdot K_i$ by
$N(D)=\prod_iN(K_i)^{d_i}$; this is well defined and multiplicative because $J_\cL$ is free.
Candidates on generators: $N(K_i)=e^{\Vol(M\setminus K_i)}$ for hyperbolic components,
$e^{2\pi\deg\Delta_{K_i}}$, or $e^{c\,\ell(\mu_i)}$.
\end{definition}

\begin{remark}[Geometric only on generators]\label{rem:normwarn}
Multiplicativity is imposed. For $D=K_1+K_2$ the value $N(K_1)N(K_2)$ is not
$e^{\Vol(M\setminus(K_1\cup K_2))}$, hyperbolic volume being neither additive over disjoint
unions nor monotone in the naive direction. The geometric reading survives on generators
only; $\Ns$-multiplicativity, a theorem about ideals, is here a definition.
\end{remark}

\begin{definition}[The system]\label{def:system}
Let
\[
G_\cL:=\varprojlim_{F,n}H_1\bigl(M\setminus\cL_F;\Z/n\Z\bigr),
\qquad
\sigma_{K_i}:=[\lambda_i]\in G_\cL,
\]
the limit over finite sublinks $F$ and integers $n$, the transition maps being induced by the
inclusions $M\setminus\cL_{F'}\subset M\setminus\cL_F$ for $F\subset F'$; $G_\cL$ is a compact
abelian group, and $[\lambda_i]$ is the compatible family of classes of the longitude of $K_i$
in the complements of the finite sublinks containing $K_i$. For finite $\cL$, $G_\cL=H_1(X;\hZ)$.
Let $X_\cL=\overline\N^{\,\cL}$ with $\overline\N=\N\cup\{\infty\}$, on which $J_\cL$ acts
by $v\mapsto v+D$, and put
\[
Y_\cL:=X_\cL\times G_\cL,\qquad D\cdot(v,g)=\bigl(v+D,\ g+\textstyle\sum_id_i\sigma_{K_i}\bigr),
\]
\[
\cA_{M,\cL}:=C(Y_\cL)\rtimes J_\cL,\qquad \sigma_t(\mu_D)=N(D)^{it}\mu_D ,
\]
the semigroup crossed product being the unital full corner of the groupoid algebra of the
dilated action of $I_\cL=\Z^{(\cL)}$, exactly as $\cA_{K,S}=C(Y_S)\rtimes J_S$ in
\cite[\S2]{Main}; measure-theoretically $Y_S=\widehat{\cO}_S\times_{U_S}G_S$ is $\N^S\times G_S$
with the same action, by \cite[Lem.~2.10]{Main}.
\end{definition}

\begin{proposition}[No new operator algebra is needed]\label{prop:paperVII}
$\cA_{M,\cL}$ is the system of \cite{Main} attached to the quadruple
$(\cL,J_\cL,G_\cL,N)$ with $\sigma_{K_i}$ in place of $\sigma_\pp$: it is the abelian case of
the Galois-twisted groupoid of Paper VII of the series with $\Frob_{K_i}=[\lambda_i]$. The
classification theorems \cite[Thm.~3.6, Thm.~3.9]{Main} and the factor-type results of
\cite{SeriesI} that depend only on the sequence $(N(K_i),\chi(\sigma_{K_i}))$ apply verbatim.
Results of the series whose proofs use arithmetic input --- Chebotarev density, prime number
theorems, class numbers, the finiteness of the unit group --- do not transfer automatically
and have to be re-proved, or fail, in the topological setting.
\end{proposition}

\begin{proof}
The proofs of \cite[Thm.~3.6, Thm.~3.9]{Main} use only that $J_S$ is a free abelian monoid on
a countable set acting on $\overline\N^{\,S}\times G_S$ as above, that $G_S$ is compact abelian
with prescribed elements $\sigma_\pp$, that $N$ is multiplicative with values in $[2,\infty)$,
and the Kolmogorov zero-one law for the forced product measure; they do not use that the
$\sigma_\pp$ generate a dense subgroup, since $\Xi_\beta$ is defined by characters of the whole
group. The quadruple $(\cL,J_\cL,G_\cL,N)$ satisfies all of this. No property of prime ideals
is used in the classification; arithmetic enters only when the criterion is evaluated.
\end{proof}

\section{The linking obstruction group}\label{sec:obstruction}

\begin{theorem}[Topological KMS classification]\label{thm:kms}
The $\mathrm{KMS}_\beta$ simplex of $\cA_{M,\cL}$ is affinely isomorphic to
$\Prob\bigl(G_\cL/\Xi_\beta(M,\cL)^\perp\bigr)$, where
\begin{equation}\label{eq:Xi}
\Xi_\beta(M,\cL)=\Bigl\{\chi\in\widehat{G_\cL}:\ \sum_iN(K_i)^{-\beta}\bigl|1-\chi([\lambda_i])\bigr|^2<\infty\Bigr\},
\end{equation}
and the state is unique if and only if every nontrivial $\chi$ has
$\sum_{i\in D_\chi}N(K_i)^{-\beta}=\infty$, where $D_\chi=\{i:\chi([\lambda_i])\ne1\}$.
\end{theorem}

\begin{proof}
Proposition~\ref{prop:paperVII} together with \cite[Thm.~3.6, Thm.~3.9]{Main}. Note that a
character trivial on every $\sigma_{K_i}$ has $D_\chi=\emptyset$ and lies in $\Xi_\beta$ for
all $\beta$; the criterion therefore fails, at every temperature, unless the $\sigma_{K_i}$
generate a dense subgroup of $G_\cL$. In the arithmetic case this density is a form of
Chebotarev's theorem; here it is a condition on the linking matrix.
\end{proof}

\begin{theorem}[The Gauss linking matrix formula]\label{thm:linking}
Let $M=S^3$, so $\widehat{G_\cL}=\bigoplus_j\Q/\Z$ by Alexander duality. For $\chi=(a_j)$ with
finite support $F$ and common denominator $m$, writing $b_j=ma_j$,
\begin{equation}\label{eq:link}
\chi\bigl(\sigma_{K_i}\bigr)=\exp\Bigl(2\pi i\sum_{j\in F}\Lk(K_i,K_j)\,a_j\Bigr),
\qquad
D_\chi=\Bigl\{i:\ \sum_{j\in F}\Lk(K_i,K_j)\,b_j\not\equiv0\ (\mathrm{mod}\ m)\Bigr\}.
\end{equation}
For $\chi$ of order $m$ the weight $|1-\chi([\lambda_i])|^2$ is bounded above and below by
positive constants on $D_\chi$, so
\[
\chi\in\Xi_\beta(S^3,\cL)\iff\sum_{i\in D_\chi}N(K_i)^{-\beta}<\infty .
\]
\end{theorem}

\begin{proof}
For $M=S^3$ and each finite sublink, $H_1(M\setminus\cL_F;\Z)$ is free on the meridians
$\mu_j$, $j\in F$, by Alexander duality, the transition maps kill the meridians of the
omitted components, and $G_\cL=\prod_j\hZ\mu_j$, whose dual is the direct sum of the duals of
the factors. Then $\iota_*\lambda_i=\sum_{j\in F}\Lk(K_i,K_j)\mu_j$ in each finite complement,
by Theorem~\ref{thm:hlhd} with our framing convention; evaluate $\chi$. The weight takes
finitely many values, all nonzero on $D_\chi$, since $\chi$ has order $m$.
\end{proof}

\begin{theorem}[Infinite detecting sets exist]\label{thm:infinite}
Let $K\subset S^3$ be a knot, $f\in\Z$ a framing, and let $\cL=\{K_1,K_2,\dots\}$ consist of
infinitely many parallel copies of $K$ with framing $f$, accumulating on $K$, so that
$\Lk(K_i,K_j)=f$ for all $i\ne j$. Let $\chi=(a_j)_{j\in F}$ and $\Sigma_F=\sum_{j\in F}a_j$.
If $f\Sigma_F\notin\Z$ then
\[
D_\chi\supseteq\{i:i\notin F\},
\]
so $D_\chi$ is cofinite and $\chi\in\Xi_\beta(S^3,\cL)$ if and only if
$\sum_iN(K_i)^{-\beta}<\infty$. Consequently $\Xi_\beta$ is genuinely $\beta$-dependent and
changes at the convergence exponent of the norm sequence. If $f\Sigma_F\in\Z$, however,
$D_\chi\subseteq F$ is finite and $\chi\in\Xi_\beta$ for every $\beta$; such characters exist
for every $f$ (take $F=\{1,2\}$, $a_1=a_2=\tfrac12$), so for this family the
$\mathrm{KMS}_\beta$ state is never unique.
\end{theorem}

\begin{proof}
For $i\notin F$ we have $\Lk(K_i,K_j)=f$ for every $j\in F$, so by \eqref{eq:link}
$\chi(\sigma_{K_i})=\exp(2\pi if\Sigma_F)$, which is $\ne1$ iff $f\Sigma_F\notin\Z$. The rest
is Theorem~\ref{thm:linking}.
\end{proof}

\begin{lemma}[Every symmetric matrix is a linking matrix]\label{lem:realize}
Let $(\ell_{ij})_{i\ne j}$ be any symmetric family of integers indexed by pairs of distinct
positive integers. There is a countable family $\cL=\{K_1,K_2,\dots\}$ of pairwise disjoint
tame knots in $S^3$, each individually unknotted, with $\Lk(K_i,K_j)=\ell_{ij}$ for all
$i\ne j$. If some column $(\ell_{ij})_i$ has infinite support, the family accumulates on
$K_j$.
\end{lemma}

\begin{proof}
Construct $K_n$ inductively, $K_1$ an unknot. Given the disjoint tame knots $K_1,\dots,K_{n-1}$,
take for each $j<n$ with $\ell_{nj}\ne0$ a collection of $|\ell_{nj}|$ small meridian circles
of $K_j$, oriented so that each links $K_j$ with linking number $\operatorname{sgn}\ell_{nj}$,
chosen of radius less than $2^{-n}$ and disjoint from everything constructed so far; if all
$\ell_{nj}=0$ take one small unknotted circle far from the previous components. Join these
circles by bands along disjoint arcs in the complement of $K_1\cup\dots\cup K_{n-1}$ and of the
circles, orienting the bands compatibly; the band-connected sum of finitely many disjoint
oriented circles is a single tame knot $K_n$, unknotted since the circles bound disjoint disks
which the bands can be chosen to avoid, and it is homologous in $S^3\setminus K_j$ to the sum
of the circles, so $\Lk(K_n,K_j)=\ell_{nj}$ for every $j<n$; by symmetry of the linking number
this fixes $\Lk(K_j,K_n)$ as well. The components are pairwise disjoint by construction, and
if column $j$ is infinitely supported, infinitely many $K_n$ pass within $2^{-n}$ of $K_j$.
\end{proof}

\begin{theorem}[A link with the Bost--Connes phase diagram]\label{thm:BC}
Let $\cL$ be a link in $S^3$ realizing, by Lemma~\ref{lem:realize}, the linking numbers
\[
\Lk(K_i,K_j)=\mathrm{bit}_{\min(i,j)}\bigl(\max(i,j)\bigr)\in\{0,1\}\qquad(i\ne j),
\]
where $\mathrm{bit}_k(n)$ is the coefficient of $2^k$ in the binary expansion of $n$, and
let $N$ be a norm with $N(K_i)=(i+1)^\gamma$ for some $\gamma>0$, so that
$\sum_iN(K_i)^{-\beta}$ converges iff $\beta>\beta_c:=1/\gamma$. Then
\[
\Xi_\beta(S^3,\cL)=\{1\}\ \text{ for }\beta\le\beta_c,\qquad
\Xi_\beta(S^3,\cL)=\widehat{G_\cL}\ \text{ for }\beta>\beta_c :
\]
the $\mathrm{KMS}_\beta$ state of $\cA_{S^3,\cL}$ is unique for $\beta\le\beta_c$, and for
$\beta>\beta_c$ the simplex is $\Prob(G_\cL)=\Prob\bigl(\prod_j\hZ\bigr)$, on which $G_\cL$
acts freely and transitively on the extreme points. Every column of the linking matrix is
infinitely supported, so $\cL$ accumulates on every one of its components.
\end{theorem}

\begin{proof}
For $\beta>\beta_c$ every sum in \eqref{eq:Xi} converges, so $\Xi_\beta=\widehat{G_\cL}$ and
$\Xi_\beta^\perp=0$. Let $\beta\le\beta_c$ and let $\chi\ne1$ have support $F$, denominator
$m\ge2$ and $b=(b_j)_{j\in F}$ not all $\equiv0\pmod m$. For $i>\max F$, \eqref{eq:link} gives
$i\in D_\chi$ iff $\sum_{j\in F}\mathrm{bit}_j(i)\,b_j\not\equiv0\pmod m$. Put $L=\max F+1$.
On any block $[k2^L,(k+1)2^L)$ of $2^L$ consecutive integers the residue $i\bmod 2^L$ runs
over all values once, so the vector $(\mathrm{bit}_j(i))_{j\in F}\in\{0,1\}^F$ takes each value
exactly $2^{L-|F|}$ times. Fix $j_0\in F$ with $b_{j_0}\not\equiv0$ and pair each
$\varepsilon\in\{0,1\}^F$ with the vector $\varepsilon'$ differing from it in the $j_0$-th
coordinate only: $\sum_j\varepsilon_jb_j$ and $\sum_j\varepsilon'_jb_j$ differ by
$\pm b_{j_0}\not\equiv0$, so at most one of them is $\equiv0$. Hence every block contains at
least $2^{L-1}$ elements of $D_\chi$, apart from the finitely many $i\le\max F$, and $D_\chi$
has lower asymptotic density at least $\frac12$. A set $D$ of positive lower density $c$ has
$k$-th element at most $k/c$ for $k$ large, hence
$\sum_{i\in D}(i+1)^{-s}\ge\sum_k(k/c+1)^{-s}=\infty$ for every $s\le1$, and $\gamma\beta\le1$;
so $\chi\notin\Xi_\beta$. Thus $\Xi_\beta=\{1\}$ and the state is unique by
Theorem~\ref{thm:kms}. Column $j$ of the matrix contains $\mathrm{bit}_j(i)$ for all $i>j$,
which is $1$ for infinitely many $i$.
\end{proof}

\begin{remark}
Theorem~\ref{thm:BC} is the topological Bost--Connes phase diagram in its original shape:
below and at the critical temperature a unique equilibrium state, above it spontaneous
symmetry breaking with the full compact group of symmetries acting simply transitively on the
extreme states, and no intermediate regime. Intermediate chains $\{\Xi_\beta\}$, as realized
arithmetically in \cite[Thm.~5.1]{SeriesI} by congruence conditions on primes, can be
realized by prescribing rows of the linking matrix whose congruence classes are controlled;
we do not pursue this. What Theorem~\ref{thm:BC} cannot control is the norm: the norm is
prescribed here, and Section~\ref{sec:volume} records what happens when it is asked to be
geometric.
\end{remark}

\begin{remark}[Why local finiteness had to go]\label{rem:whyaccum}
If $\cL$ were locally finite, a compact Seifert surface $\Sigma_j$ for $K_j$ would meet only
finitely many components, so $\Lk(K_j,K_i)\ne0$ for finitely many $i$; by symmetry of the
linking number every column of the matrix would then be finitely supported and every $D_\chi$
finite, returning the empty theory of \S\ref{sec:autopsy}. By Remark~\ref{rem:accumulate}
local finiteness is in any case impossible for infinite $\cL$ in closed $M$. Accumulation is
what lets $\Sigma_j$ be pierced infinitely often, and Theorem~\ref{thm:infinite} realizes
this in the simplest way; Lemma~\ref{lem:realize} shows that nothing beyond symmetry
constrains the linking matrix of a countable tame link.
\end{remark}

\begin{remark}[Numerical illustration]\label{rem:numdensity}
For random integer linking data $\Lk(K_i,K_j)\in[-4,4]$ over $400$ components and a character
supported on two meridians, the detecting set contained $45$--$47\%$ of components for $m=2$,
$65$--$69\%$ for $m=3$ and $78$--$90\%$ for $m=5$. This illustrates the generic behaviour of
the congruence condition in \eqref{eq:link}; it is not itself a link realization, the random
matrices being constrained neither to symmetry nor to arising from an embedded family. The
rigorous statements are Theorems~\ref{thm:infinite} and~\ref{thm:BC}.
\end{remark}

\section{Ambient homological obstructions}\label{sec:ambient}

Let $\nu:G_\cL\to H_1(M;\hZ)$ be the map induced by the inclusions $M\setminus\cL_F\subset M$;
it is surjective, since $H_1(M\setminus\cL_F)\to H_1(M)$ is onto for every finite $F$ and the
groups are compact, and $\nu(\sigma_{K_i})=[K_i]$, $\nu([\mu_i])=0$. As in
Theorem~\ref{thm:dict} put $\Cl_\cL=\langle[K_i]:i\rangle\le H_1(M;\Z)$.

\begin{theorem}[Ambient homology]\label{thm:ambient}
Let $\psi$ be a character of $H_1(M;\hZ)$ and $\chi=\psi\circ\nu\in\widehat{G_\cL}$. Then
$\chi(\sigma_{K_i})=\psi([K_i])$ for every $i$, and
\[
\chi\in\Xi_\beta(M,\cL)\iff\sum_{i:\ \psi([K_i])\ne1}N(K_i)^{-\beta}<\infty .
\]
Consequently:
\begin{enumerate}
\item If $\psi$ is trivial on $\Cl_\cL$, then $\chi\in\Xi_\beta(M,\cL)$ for every $\beta>0$.
The characters of $H_1(M;\hZ)/\overline{\Cl_\cL}$ pull back injectively to $\Xi_\beta$, so
with $h^0_\cL:=|H_1(M;\Z)/\Cl_\cL|$ the $\mathrm{KMS}_\beta$ simplex has at least $h^0_\cL$
extreme points at every temperature, and uniqueness never holds unless the knot classes
generate $H_1(M;\Z)$. This obstruction has no arithmetic counterpart: $G_S$ maps onto the
subgroup $\Cl^+_S$ generated by the primes of $S$ and never sees $\Cl^+_K/\Cl^+_S$.
\item If $\psi$ is nontrivial on $\Cl_\cL$ and the knots whose classes lie outside
$\ker\psi\cap\Cl_\cL$ form a $\beta$-summable set, then $\chi\in\Xi_\beta(M,\cL)$ and the
$\mathrm{KMS}_\beta$ state is not unique. This is the exact counterpart of the unramified
class-group obstruction \cite[Prop.~4.1]{Main}: a proper subgroup $H<\Cl_\cL$ containing the
classes of all but a $\beta$-summable set of components, with $\Cl_\cL$ in place of $\Cl^+_S$
and $[K_i]$ in place of $[\pp]$.
\end{enumerate}
\end{theorem}

\begin{proof}
$\chi(\sigma_{K_i})=\psi(\nu[\lambda_i])=\psi([K_i])$. The group $H_1(M;\hZ)$ is profinite,
so $\psi$ has finite order and $|1-\psi([K_i])|^2$ takes finitely many values, bounded below
by a positive constant where it is nonzero; the displayed equivalence follows from
\eqref{eq:Xi}. (1): all summands vanish. Distinct characters of
$H_1(M;\hZ)$ pull back to distinct characters of $G_\cL$ since $\nu$ is surjective, and a
group of $h^0_\cL$ characters inside $\Xi_\beta$ forces $|G_\cL/\Xi_\beta^\perp|\ge h^0_\cL$.
(2): the sum in the equivalence runs over the components with $[K_i]\notin\ker\psi$.
\end{proof}

\begin{corollary}[The parallel, corrected]\label{cor:parallel}
The counterpart of the unramified obstruction of \cite[\S4.1]{Main} is
Theorem~\ref{thm:ambient}(2), under
\[
\Cl^+_S=\bigl\langle[\pp]:\pp\in S\bigr\rangle\le\Cl^+_K
\ \longleftrightarrow\
\Cl_\cL=\bigl\langle[K_i]:i\bigr\rangle\le H_1(M;\Z),
\qquad h^+_S=|\Cl^+_S|\leftrightarrow|\Cl_\cL| .
\]
The quotient $H_1(M;\Z)/\Cl_\cL$ of Theorem~\ref{thm:ambient}(1) is a purely topological
addition: the complement of a link remembers the homology of the ambient manifold, while the
$S$-ideles remember nothing of the class group beyond the classes of $S$.
\end{corollary}

\begin{example}\label{ex:lens}
Let $M=L(p,1)$, $H_1(M;\Z)=\Z/p$. If all $K_i$ are null-homologous, $\Cl_\cL=0$, $h^0_\cL=p$,
and the simplex has at least $p$ extreme points at every $\beta$ by
Theorem~\ref{thm:ambient}(1); the $h^+_S$-type obstruction is absent, as for a field with
$\Cl^+_S=1$. If instead some components generate $H_1(M)$, then $h^0_\cL=1$, $\Cl_\cL=\Z/p$,
and Theorem~\ref{thm:ambient}(2) applies: for a character $\psi$ of $\Z/p$, the pull-back
lies in $\Xi_\beta$ iff the components with $\psi([K_i])\ne1$ are $\beta$-summable, so a link
in which all but a summable family of components are null-homologous has broken symmetry at
every $\beta$, exactly as an $S$ whose primes lie, up to a summable set, in a proper subgroup
of $\Cl^+_S$. What matters is the distribution of the homology classes of the knots, as it
is the distribution of the ideal classes of the primes. For $M=S^3$ both obstructions are
absent, and by Theorems~\ref{thm:infinite} and~\ref{thm:BC} the whole phase structure comes
from linking.
\end{example}

\section{Volume exponents and closed phases}\label{sec:volume}

\begin{definition}\label{def:volexp}
For hyperbolic $\cL$ set
$\beta_c(\cL)=\inf\{\beta>0:\sum_ie^{-\beta\Vol(M\setminus K_i)}<\infty\}$.
\end{definition}

\begin{proposition}[The full family is too large]\label{prop:JT}
Over the family of all hyperbolic knots in $S^3$ the series $\sum_ie^{-\beta\Vol_i}$
diverges for every $\beta>0$, so $\beta_c=\infty$, and a finite volume exponent requires a
sparse subfamily.
\end{proposition}

\begin{proof}
By Thurston's hyperbolic Dehn surgery theorem \cite[Ch.~5]{Thurston} the twist knots $K_n$,
obtained by $1/n$-surgery on one component of the Whitehead link, are hyperbolic for $n$
large with volumes increasing to the volume $V_W\approx3.66$ of the Whitehead link
complement. So infinitely many hyperbolic knots have volume below $V_W$, and the series has
infinitely many terms bounded below by $e^{-\beta V_W}$.
\end{proof}

\begin{remark}
The count of hyperbolic $3$-manifolds of volume at most $V$ is not finite in dimension $3$,
precisely because of Dehn filling; the growth rate $V^{cV}$ of Burger--Gelander--Lubotzky--Mozes
concerns dimension $\ge4$. The divergence in Proposition~\ref{prop:JT} is therefore stronger
than a growth estimate: the volume spectrum of knots has accumulation points, and the volume
series diverges term by term near each of them.
\end{remark}

\begin{proposition}[Open and closed broken phases]\label{prop:closed}
Let $\cL$ be the link of Theorem~\ref{thm:BC} and let $N$ be any norm with
$\sum_{i\in D}N(K_i)^{-\beta}=\infty$ for every $\beta<\beta_c$ and every set $D$ of positive
lower density, where $\beta_c$ is the convergence exponent of $\sum_iN(K_i)^{-\beta}$. Then
the $\mathrm{KMS}_\beta$ state is unique for $\beta<\beta_c$ and the simplex is $\Prob(G_\cL)$
for $\beta>\beta_c$, and at $\beta=\beta_c$ the state is unique if $\sum_iN(K_i)^{-\beta_c}=\infty$
and the simplex is $\Prob(G_\cL)$ if $\sum_iN(K_i)^{-\beta_c}<\infty$: the broken phase is
the open half-line $(\beta_c,\infty)$ for $N(K_i)=(i+1)^{1/\beta_c}$ and the closed half-line
$[\beta_c,\infty)$ for $N(K_i)=\bigl((i+1)\log^2(i+2)\bigr)^{1/\beta_c}$, exactly as in
\cite[Thm.~5.2]{Main}.
\end{proposition}

\begin{proof}
The proof of Theorem~\ref{thm:BC} shows that every nontrivial character is detected on a set
of lower density at least $\frac12$, so for $\beta<\beta_c$ the hypothesis on $N$ gives
$\Xi_\beta=\{1\}$; for $\beta>\beta_c$ all sums converge. At $\beta_c$: if the full series
diverges, so does every sum over a set $D$ of positive lower density $c$ for the two norms
named, since their terms $a_i$ are eventually decreasing and the $k$-th element of $D$ is at
most $k/c$, so that $\sum_{i\in D}a_i\ge\sum_ka_{\lceil k/c\rceil}\ge\int_1^\infty a(x/c)\,dx=c\int_{1/c}^\infty a=\infty$;
if it converges, all subsums converge. The two norms have
$\sum_i(i+1)^{-1}=\infty$ and $\sum_i\bigl((i+1)\log^2(i+2)\bigr)^{-1}<\infty$, and for
$\beta<\beta_c$ their $\beta$-sums over positive-density sets diverge by partial summation.
\end{proof}

\begin{remark}\label{rem:volopen}
With a prescribed norm the theory is complete: Theorem~\ref{thm:BC} and
Proposition~\ref{prop:closed} realize both the open and the closed phase diagram of
\cite[Thm.~5.2]{Main} by linking alone. What we have not exhibited is a family that is
simultaneously linking-rich in the sense of Theorem~\ref{thm:BC} and \emph{geometric} in its
norm: the components of Lemma~\ref{lem:realize} are unknots, and controlling the hyperbolic
volumes of an accumulating family --- whose complements degenerate --- is the concrete gap on
the geometric side. Definition~\ref{def:norm} allows any norm, but the dictionary's promise
that $\Ns\pp$ corresponds to a geometric quantity is kept only on generators, by
Remark~\ref{rem:normwarn}, and not yet by any accumulating family.
\end{remark}

\section{Contrastive autopsy: why the meridian fails}\label{sec:autopsy}

\begin{theorem}[The inertia model has no thermodynamics]\label{thm:meridian}
Let $M=S^3$ and define the system with $\sigma_{K_i}=[\mu_i]$ instead. Then every
$\chi\in\widehat{G_\cL}$ has finite detecting set, $\Xi_\beta=\widehat{G_\cL}$ for every
$\beta>0$, the $\mathrm{KMS}_\beta$ state is unique at every temperature, and no phase
transition occurs --- for any link, any norm, any $\beta$.
\end{theorem}

\begin{proof}
As in the proof of Theorem~\ref{thm:linking}, $G_\cL\cong\prod_i\hZ\mu_i$ and
$\widehat{G_\cL}=\bigoplus_i\Q/\Z$. A character is thus nontrivial in finitely many
coordinates, and since $\mu_i$ is the $i$-th basis vector, $D_\chi$ is that finite set.
\end{proof}

\begin{proposition}[The arithmetic contrast]\label{prop:spread}
Let $K=\Q$. In $G_S=\Idl_S/\overline{\Q^\times_{S,+}}$, with the convention
$\sigma_\aaa=r(s_\aaa)^{-1}$ of \cite[(2.2)]{Main},
\begin{equation}\label{eq:sigmap}
\sigma_p=\theta_p=\bigl(p\bigr)_{\ell\ne p}\times(1)_{\ell=p}\in U_S=\prod_\ell\Z_\ell^\times ,
\end{equation}
nontrivial at every place $\ell\ne p$. Consequently a character supported at a single place
$\ell_0$, with component $\chi_{\ell_0}$ a character of $\Z_{\ell_0}^\times$ of finite order,
has $\chi(\sigma_p)=\chi_{\ell_0}(p)=\chi(\Frob_p)^{-1}$ for $p\ne\ell_0$, with detecting set
of positive Dirichlet density.
\end{proposition}

\begin{proof}
$\sigma_p$ is the class of the inverse of the idele $s_p$ with $p$ at $p$ and $1$ elsewhere;
since $p$ lies in $\Q^\times_{S,+}$ the diagonal idele $(p,p,\dots)$ is trivial in $G_S$, and
$p^{-1}s_p$ has component $p^{-1}$ at $\ell\ne p$ and $1$ at $p$, so its inverse is
\eqref{eq:sigmap}; this is \cite[Lem.~2.10(1)]{Main}. Evaluating $\chi$ at \eqref{eq:sigmap}
gives $\chi_{\ell_0}(p)$, and $\chi(\sigma_p)=\chi(\Frob_p)^{-1}$ is \cite[Lem.~3.8]{Main}.
A nontrivial $\chi_{\ell_0}$ factors through $(\Z/\ell_0^k)^\times$, so $D_\chi$ is a union
of arithmetic progressions, of positive Dirichlet density by Dirichlet's theorem
\cite[Ch.~VII]{Neukirch}.
\end{proof}

\begin{remark}[The moral]\label{rem:moral}
The two computations are the same statement. Arithmetically it is the product formula that
spreads a one-place idele across all others; topologically it is the relation
$\lambda_i=\sum_j\Lk(K_i,K_j)\mu_j$ that spreads the longitude across all linked components.
\emph{Inertia is local and carries no thermodynamics; Frobenius is global and carries the
linking matrix.} The meridian model of Theorem~\ref{thm:meridian} is not so much a failed
construction as a correct computation of the wrong object: it says inertia classes are
thermodynamically inert, which is what inertia is for. Numerically, for
$\ell_0=5,7,11$ with $\chi_{\ell_0}$ the Legendre symbol the arithmetic detecting sets
contained $51.5\%$, $52.8\%$ and $52.2\%$ of the primes below $2000$, matching the
topological densities of Remark~\ref{rem:numdensity}.
\end{remark}

\section{The dictionary, and assessment}\label{sec:dict}

\[
\begin{array}{@{}l@{\ \ }l@{\ \ }l@{}}
\text{Arithmetic \cite{Main}} & \text{Knot topology} & \text{Reference}\\\hline
\pp\in S & K_i\in\cL & \text{---}\\
K_\pp^\times & H_1(T_i;\hZ) & \text{Lem.~1.1}\\
\cO_\pp^\times\ (\text{inertia}) & \mu_i\ (\text{meridian}) & \text{Lem.~1.1}\\
\pi_\pp^{\hZ}\ (\text{Frobenius}) & \lambda_i\ (\text{longitude}) & \text{Lem.~1.1}\\
\Idl_S & H_1(\partial X;\hZ) & \text{Thm.~1.4}\\
\overline{K^\times_{S,+}} & \ker\iota_*=\langle\tlam_i\rangle & \text{Thm.~1.3, 1.4}\\
\overline{\cO^\times_{K,+}} & \ker\iota_*\cap\bigoplus\hZ\mu_i;\ =0\text{ for }S^3 & \text{Thm.~1.4}\\
\Cl^+_S=\langle[\pp]\rangle\le\Cl^+_K & \Cl_\cL=\langle[K_i]\rangle\le H_1(M;\Z) & \text{Thm.~1.4, 4.1}\\
h^+_S>1\ \text{obstruction} & [K_i]\text{ summably in a proper }H<\Cl_\cL & \text{Thm.~4.1(2)}\\
\text{---} & H_1(M;\Z)/\Cl_\cL\ne0:\ \text{broken at all }\beta & \text{Thm.~4.1(1)}\\
\sigma_\pp=(\pi_\pp)_{\qq\ne\pp} & \lambda_i=\sum_j\Lk(K_i,K_j)\mu_j & \text{Prop.~6.2, Thm.~3.2}\\
\text{product formula} & \text{linking relation} & \text{Rem.~6.3}\\
\zeta_{K,S}(\beta) & \sum_iN(K_i)^{-\beta},\ N\text{ geometric on generators} & \text{Def.~2.2, 5.1}\\
\Xi_\beta\ \text{Dirichlet criterion} & \Xi_\beta\ \text{linking criterion} & \text{Thm.~3.1, 3.2}\\
\text{Chebotarev density} & \text{bit matrix: each }\chi\ne1\text{ detected, density}\ge\tfrac12 & \text{Thm.~\ref{thm:BC}}\\
\text{BC phase diagram} & \text{realized by linking, open or closed} & \text{Thm.~\ref{thm:BC}, Prop.~\ref{prop:closed}}
\end{array}
\]

The transfer is complete at the level of structure and, by Proposition~\ref{prop:paperVII},
requires no new operator algebra: the classification of \cite{Main} is a theorem about a
quadruple with no arithmetic in it. What the topological instantiation contributes is a
different mechanism for the same phenomenon --- linking in place of reciprocity --- and
\S\ref{sec:autopsy} shows the two are structurally identical, both being ways of spreading a
local generator over infinitely many places.

Three questions. First, realization: \cite[Thm.~5.1]{SeriesI} realizes prescribed chains
$\{\Xi_\beta\}$ from congruence conditions on primes, using Chebotarev to select primes in
prescribed classes; Lemma~\ref{lem:realize} removes every constraint on the linking matrix,
so the topological realization problem is purely combinatorial, and Theorem~\ref{thm:BC}
settles the extreme case. Second, geometry: Remark~\ref{rem:volopen} records that we have no
family which is simultaneously linking-rich and geometric in its norm, and this is where a
genuinely three-dimensional input would enter. Third, and we state it as a caution rather than a
programme: nothing here uses a quantum invariant. The Jones polynomial, Chern--Simons theory
and the volume conjecture play no role in any proof above, and the subfactors of Paper II
of the series arise from arithmetic symmetry breaking rather than from knot theory. We
see no evidence that the two occurrences of subfactors are related, and would not want the
coincidence of names to suggest otherwise.

\begin{thebibliography}{9}
\bibitem{Kapranov} M. Kapranov, A. Smirnov, \emph{Cohomology determinants and reciprocity laws}, preprint, 1995.
\bibitem{Mazur} B. Mazur, \emph{Remarks on the Alexander polynomial}, unpublished notes, 1963--64.
\bibitem{Morishita} M. Morishita, \emph{Knots and Primes: An Introduction to Arithmetic Topology}, Universitext, Springer, 2012.
\bibitem{Neukirch} J. Neukirch, \emph{Algebraic Number Theory}, Grundlehren 322, Springer, 1999.
\bibitem{Reznikov} A. Reznikov, \emph{Three-manifolds class field theory}, Selecta Math. (N.S.) 3 (1997), 361--399.
\bibitem{Thurston} W. P. Thurston, \emph{The Geometry and Topology of Three-Manifolds}, Princeton lecture notes, 1980.
\bibitem{Main} R. Chen, \emph{KMS state uniqueness and phase transitions in localized Bost--Connes systems: from Chebotarev density to the Hardy--Littlewood conjecture}, preprint, 2026, DOI \href{https://doi.org/10.5281/zenodo.22152101}{10.5281/zenodo.22152101}.
\bibitem{SeriesI} R. Chen, \emph{Localized Bost--Connes systems I: factor type and the spectrum of transitions}, preprint, 2026, DOI \href{https://doi.org/10.5281/zenodo.22160827}{10.5281/zenodo.22160827}.
\end{thebibliography}

\end{document}
